\chapter{Health Checks for Women}
\index{Health Checks for Women}

\section*{Definition and Overview}
Health checks for women are routine medical assessments that address women's health across the lifespan, including reproductive health, screening for cancers that affect women, cardiovascular risk, bone health, and mental health. Key components include cervical screening, breast awareness and breast screening, discussion of contraception and family planning, pregnancy care where relevant, and assessment of menopausal health. Coverage and access to women's health services vary by health system. Regular checks detect problems early and provide the opportunity for prevention, particularly for conditions such as cervical and breast cancer, osteoporosis, and heart disease.

\section*{Epidemiology}
Women's health issues are common. Around 1 in 11 women will develop breast cancer, and around 1 in 50 will develop cervical cancer without screening; both are far more survivable when detected early. Heart disease is the leading cause of death in women. Osteoporosis affects around 1 in 4 women over 50, and menopause-related symptoms affect most women at some stage. Endometriosis affects around 1 in 10 women. Many of these conditions are asymptomatic early, which is why regular checks and screening matter.

\section*{Aetiology and Pathophysiology}
Health checks for women target the main preventable and detectable conditions.

\begin{itemize}
    \item \textbf{Cervical cancer:} caused by persistent high-risk HPV infection; detected by cervical screening and prevented by vaccination
    \item \textbf{Breast cancer:} risk increases with age; detected early by mammography screening and breast awareness
    \item \textbf{Cardiovascular disease:} risk factors include blood pressure, cholesterol, diabetes, smoking, and menopause
    \item \textbf{Osteoporosis:} bone loss accelerates after menopause; detected by bone density testing
    \item \textbf{Reproductive health:} menstrual disorders, contraception, fertility, and pregnancy
    \item \textbf{Menopause:} hormonal changes and their effects on health
    \item \textbf{Mental health:} depression and anxiety are common, including around pregnancy and menopause
    \item \textbf{Other:} thyroid disease, iron deficiency, and urinary incontinence are common in women
\end{itemize}

\section*{Clinical Features}
A women's health check covers a broad range of topics.

\begin{itemize}
    \item Cervical screening every five years from age 25 to 74
    \item Breast awareness, and mammography screening every two years from age 40--50 depending on risk
    \item Blood pressure, cholesterol, and blood glucose checks
    \item Bone density testing for women at risk of osteoporosis
    \item Discussion of contraception and family planning
    \item Menstrual and gynaecological health, including heavy, painful, or irregular periods
    \item Menopause symptoms and management options
    \item Mental health and wellbeing
    \item Vaccination review, including HPV, flu, and other vaccines
    \item Lifestyle assessment: smoking, alcohol, diet, and physical activity
\end{itemize}

\section*{Differential Diagnosis}
Findings at a health check may lead to further assessment.

\begin{itemize}
    \item Abnormal cervical screening results requiring follow-up
    \item Breast lumps or changes requiring imaging
    \item Elevated blood pressure, cholesterol, or glucose requiring confirmation
    \item Low bone density requiring osteoporosis treatment
    \item Menstrual problems requiring investigation, including endometriosis and fibroids
    \item Menopause symptoms that overlap with thyroid disease and anaemia
    \item Mental health concerns requiring assessment
    \item Over-screening and under-screening, both of which are best avoided
\end{itemize}

\section*{When to Seek Medical Care}
All women should have regular health checks, with frequency based on age and risk. See a doctor promptly for any new or concerning symptom, including breast lumps, abnormal bleeding, pelvic pain, or chest pain.

\textbf{Contact your local emergency services for:}
\begin{itemize}
    \item Chest pain, breathlessness, or symptoms of stroke
    \item Heavy vaginal bleeding with faintness or dizziness
    \item Sudden severe pelvic or abdominal pain
    \item Bleeding in pregnancy with pain or collapse
\end{itemize}

\section*{Diagnosis and Investigations}
A women's health check involves targeted history, examination, and tests.

\begin{itemize}
    \item \textbf{History:} menstrual, gynaecological, obstetric, sexual, and family history
    \item \textbf{Blood pressure, weight, and BMI:} cardiovascular risk assessment
    \item \textbf{Blood tests:} cholesterol, glucose, iron studies, thyroid function, and other tests as indicated
    \item \textbf{Cervical screening:} an HPV test every five years for women aged 25--74
    \item \textbf{Breast screening:} mammography according to the national program and personal risk
    \item \textbf{Bone density testing:} for women at risk of osteoporosis
    \item \textbf{Pelvic examination:} when symptoms or screening require it
    \item \textbf{STI screening:} based on sexual history
\end{itemize}

\section*{Management}
Management turns the findings of the check into a plan.

\begin{itemize}
    \item \textbf{Cancer prevention:} HPV vaccination, cervical screening, breast screening, and acting on results
    \item \textbf{Cardiovascular risk:} lifestyle measures and medicines for blood pressure, cholesterol, and glucose
    \item \textbf{Bone health:} calcium, vitamin D, exercise, and treatment for osteoporosis
    \item \textbf{Reproductive health:} contraception, fertility discussions, and treatment of menstrual disorders
    \item \textbf{Menopause:} symptom management, including hormone therapy where appropriate
    \item \textbf{Mental health:} treatment and referral as needed
    \item \textbf{Lifestyle:} smoking cessation, healthy eating, and physical activity
    \item \textbf{Scheduled review:} a plan for the next check and follow-up
\end{itemize}

\section*{Complications}
\begin{itemize}
    \item Missed early diagnosis of breast or cervical cancer without screening
    \item Undiagnosed cardiovascular risk, which is the leading cause of death in women
    \item Osteoporosis and fractures in untreated bone loss
    \item Delayed diagnosis of gynaecological conditions, including endometriosis
    \item Anxiety from results that are not explained
    \item Over-testing and unnecessary intervention
\end{itemize}

\section*{Prognosis and Outcomes}
Health checks for women save lives and improve health. Cervical screening and vaccination have dramatically reduced cervical cancer rates, and breast screening reduces deaths from breast cancer. Early treatment of blood pressure and cholesterol prevents heart attacks and strokes, and osteoporosis treatment prevents fractures. Women who attend regular checks and act on results have better health outcomes across the lifespan, from reproductive years through to older age.

\section*{Prevention and Self-Care}
\begin{itemize}
    \item Book a women's health check with your GP
    \item Attend cervical screening every five years from age 25
    \item Take part in breast screening from age 40--50 according to the schedule
    \item Know your family history of breast, ovarian, and bowel cancer
    \item Monitor blood pressure and cholesterol
    \item Maintain a healthy weight, exercise, and do not smoke
    \item Ensure adequate calcium and vitamin D for bone health
    \item See a doctor for abnormal bleeding, pain, or breast changes
    \item Attend follow-up and specialist appointments as advised
\end{itemize}

\section*{Sources and Further Reading}
The following reputable sources provide further information for healthcare students and patients seeking reliable, current advice about health checks for women.

\begin{itemize}
    \item Jean Hailes for Women's Health. \textit{Women's health checks and wellbeing.} https://www.jeanhailes.org.au/
    \item Australian Government Department of Health and Aged Care. \textit{Cervical and breast screening programs.} https://www.health.gov.au/
    \item Healthdirect Australia. \textit{Health checks for women.} https://www.healthdirect.gov.au/
    \item Royal Australian College of General Practitioners. \textit{Women's health guidelines.} https://www.racgp.org.au/
    \item NHS. \textit{Women's health checks.} https://www.nhs.uk/
    \item World Health Organization. \textit{Women's health.} https://www.who.int/
\end{itemize}
